\section{Conclusiones parciales}
\label{sec:conclusiones}

El trabajo realizado hasta la fecha permite extraer las siguientes conclusiones
parciales, entendidas como estado de avance y no como cierre del proyecto.

\textbf{Sobre el marco te\'orico.} Se complet\'o el estudio sistem\'atico de los m\'etodos de
ondas superficiales, con los ocho cap\'itulos del texto de referencia resumidos, m\'as de
140 conceptos at\'omicos interconectados y una base bibliogr\'afica de 65 art\'iculos
analizados individualmente. El relevamiento confirma que la teor\'ia MASW, las buenas
pr\'acticas de adquisici\'on y las herramientas de procesamiento son maduras y accesibles,
y que la barrera de entrada real para su aplicaci\'on local es el \textbf{instrumental de
adquisici\'on multicanal sincronizado}. Esto valida el enfoque del proyecto: el aporte
est\'a en el sistema de adquisici\'on, no en el m\'etodo de an\'alisis.

\textbf{Sobre el desarrollo instrumental.} Se construy\'o un sistema completo y funcional
que cubre la cadena entera, desde el transductor hasta el perfil $V_S(z)$: cadena
anal\'ogica de cuatro etapas con ganancia programable, adquisici\'on digital determinista
gobernada por una m\'aquina de estados en hardware, jerarqu\'ia de almacenamiento de tres
niveles, red inal\'ambrica multi-nodo con disparo sincronizado por flanco f\'isico,
interfaz web de operador que elimina la necesidad de una computadora en campo, y un
\emph{pipeline} de procesamiento con inversi\'on multimodal. El sistema fue validado
funcionalmente mediante una auditor\'ia sistem\'atica con criterios de integridad
verificables.

\textbf{Sobre el aporte metodol\'ogico.} El desarrollo del \emph{front-end} anal\'ogico
asistido digitalmente constituye el resultado m\'as maduro del proyecto. Aplicar la
filosof\'ia de \emph{digitally assisted analog design}~\cite{Murmann2004} ---concebida
originalmente a nivel de circuito integrado--- al nivel de instrumentaci\'on de bajo costo
result\'o efectivo: el \emph{offset} residual a la entrada del ADC se redujo de hasta
\SI{45.8}{\percent} a menos de \SI{3.9}{\percent} del fondo de escala, con una \'unica
configuraci\'on de controlador para todas las placas y ganancias evaluadas, y con
significancia estad\'istica en las ocho configuraciones. El car\'acter \emph{foreground} de
la calibraci\'on ---activa antes de adquirir y desactivada durante la adquisici\'on---
preserva un camino de se\~nal puramente anal\'ogico, evitando el compromiso entre tiempo de
establecimiento y ancho de banda propio de los servos de continua continuos. Este
resultado fue sometido a revisi\'on externa mediante su env\'io al congreso URUCON~2026.

\textbf{Sobre la metodolog\'ia de trabajo.} Dos pr\'acticas resultaron decisivas y merecen
destacarse. La primera es la \textbf{verificaci\'on automatizada} como condici\'on de
avance: las suites de pruebas de firmware y el \emph{gate} de verificaciones del software
detectaron defectos que de otro modo habr\'ian llegado a campo. La segunda es la adopci\'on
de \textbf{criterios de aceptaci\'on verificables} en lugar de inspecci\'on visual: la
comprobaci\'on de integridad por suma criptogr\'afica permiti\'o distinguir una transferencia
completa de una truncada, distinci\'on que result\'o decisiva en la auditor\'ia. Los defectos
cr\'iticos encontrados ---en particular la carrera de \SI{10}{\milli\second} que dejaba
inutilizable un nodo de forma determinista--- ilustran por qu\'e un sistema distribuido de
adquisici\'on sincronizada requiere validaci\'on sistem\'atica y no meramente funcional.

\textbf{Sobre lo pendiente.} El proyecto no est\'a cerrado. La validaci\'on que constituye
su criterio de \'exito \'ultimo ---el contraste cuantitativo de los perfiles $V_S$ obtenidos
contra un ensayo geot\'ecnico de referencia--- a\'un no se realiz\'o. Tampoco est\'an
resueltos el hardware definitivo en PCB propia ni la fuente s\'ismica repetible, ambos
frentes iniciados pero no concluidos. La caracterizaci\'on de apareamiento entre canales y
de \emph{jitter} sobre el arreglo completo, que es cr\'itica para la calidad de la imagen
de dispersi\'on, est\'a hasta ahora acotada s\'olo por presupuesto te\'orico y por medici\'on
sobre pocos nodos.

\textbf{Balance.} El estado de avance corresponde a un sistema que \textbf{funciona de
extremo a extremo y produce resultados}, con una parte de la instrumentaci\'on validada
cuantitativamente y sometida a revisi\'on externa, y con los frentes restantes
identificados, acotados y ya iniciados. El trabajo restante es de consolidaci\'on y
validaci\'on externa, no de exploraci\'on.
